\chapter{Maps: Scale as a Declared Discard Set}
\label{ch:maps}

\begin{quotation}
\noindent\textit{Synopsis.} Maps are examined here as one of the oldest and
most successful examples of explicit, declared information reduction. A
map's scale functions as a declared discard set, informing the user in
advance which distinctions have been intentionally omitted. The chapter
shows that different scales support different continuation tasks, that
insufficiency is not equivalent to error, and that a map can be perfectly
adequate for one purpose while being wholly inadequate for another ---
underscoring how transparency alone can turn what would otherwise be
corrupting loss into benign simplification.
\end{quotation}

\section{The Scale Bar as a Prehistoric Loss Ledger}

\begin{definition}[Scale-Induced Discard]
\label{def:scale-discard}
A map at scale $1{:}s$ discards every distinction $d$ whose minimal pair
(\cref{def:minimal-pair}) corresponds to a real-world separation below a
resolution threshold $\theta(s)$, where $\theta$ is increasing in $s$.
\end{definition}

\begin{proposition}[Scale as Declared Discard]
\label{prop:scale-declared}
A map's scale declaration is an instance of the loss ledger
(\cref{def:loss-ledger}).
\end{proposition}

\begin{proofsketch}
By \cref{def:scale-discard}, a stated scale specifies, implicitly, the
class of spatial distinctions below threshold $\theta(s)$ that the map
does not represent. This is functionally equivalent to a declared
discard set (\cref{def:discard-set}) attached openly to the rendering,
satisfying \cref{def:benign-loss} rather than \cref{def:corrupting-loss}
for any task not requiring sub-threshold detail.
\end{proofsketch}

\section{Task-Relative Sufficiency, Concretely}

\begin{example}
A $1{:}50{,}000$ map has $\theta(50{,}000)$ on the order of tens of
meters: two buildings on the same block collapse to the same mark, but
two distinct towns do not. This map is sufficient
(\cref{def:sufficiency}) for route planning between towns and
insufficient for locating a specific building on a specific block,
without being incorrect at either task -- the discarded distinctions are
declared by the scale bar itself, and \cref{prop:corruption-undeclared}
confirms this declared loss does not constitute corruption regardless of
which task a given user brings to the map. A $1{:}5{,}000$ map lowers
$\theta$ by an order of magnitude, becoming sufficient for the
building-level task at the cost of no longer fitting an entire province
on one page -- a different, but equally declared, discard set.
\end{example}

\section{Completing the Taxonomy}

Chapter~\ref{ch:painting} previewed a two-by-two structure crossing
closed against open hidden state with declared against undeclared
discard, and placed maps in the closed-and-declared cell. It is worth
saying precisely why maps belong there and not in the open cell painting
occupies, since terrain is, considered in the abstract, an enormously
complex physical domain -- arguably as complex as anything a painting
depicts.

\begin{proposition}[Map Closedness Is a Property of the Discard Function, Not the Domain]
\label{prop:map-closedness}
A map's hidden-state space is closed (\cref{def:closed-open}) regardless
of the physical complexity of the terrain it depicts.
\end{proposition}

\begin{proofsketch}
By \cref{def:scale-discard}, everything a map discards is characterized
by a single, known threshold $\theta(s)$ applied uniformly: the discarded
distinction set is exactly $\{d : d\text{'s separation} < \theta(s)\}$,
which is a well-defined, if large, set fixed by $s$ alone. This is an
enumerable characterization of what is missing, satisfying
\cref{def:closed-open}'s closed case, regardless of how much real-world
detail falls below $\theta(s)$. Painting's causal-history space
(\cref{prop:open-causal-space}) is open not because it is \emph{large}
but because no threshold-based or otherwise finite characterization of
the missing causes exists at all; terrain's complexity, in contrast, is
fully characterized by $\theta(s)$ once the scale is fixed.
\end{proofsketch}

\cref{prop:map-closedness} is the precise statement behind
Chapter~\ref{ch:painting}'s claim that closedness is a property of the
\emph{discard declaration}, not of the depicted domain's complexity: a
map's terrain could in principle be as intricate as a painted scene, and
the map would still be closed, because what makes a domain open in this
monograph's sense is the absence of any finite characterization of what
is missing, not the sheer quantity of what is missing.
